\section*{Streszczenie}
Praca opisuje pomiar wpływu trzech pozycji głowy oraz czterech długości ślizgu po odbiciu od ściany na czas przepłynięcia 25 m stylem dowolnym. Pomiary wykonał jeden zawodnik, autor pracy. Najkrótszy średni czas uzyskano przy neutralnej pozycji głowy i ślizgu o długości 4 m. Zawodnik przed badaniem stosował pozycję z głową uniesioną, która okazała się najwolniejsza. Praca została zatwierdzona przez dr. Roberta Żukawskiego.

\textbf{Słowa kluczowe:} pływanie, styl dowolny, pozycja ciała, opór hydrodynamiczny, ślizg.

\section*{1. Wstęp}
Prędkość pływaka zależy od siły napędowej wytwarzanej przez ruchy rąk i nóg oraz od oporu wody. Opór dzieli się na trzy składowe. Opór kształtu wynika z różnicy ciśnień przed ciałem i za nim. Opór falowy powstaje na granicy wody i powietrza. Opór tarcia wynika z przepływu wzdłuż powierzchni ciała. Pozycja ciała wpływa głównie na opór kształtu. Zmienia ona pole przekroju ciała prostopadłego do kierunku ruchu.

W stylu dowolnym za korzystną uznaje się pozycję wyprostowaną i poziomą. Ciało ma tworzyć jedną linię od głowy do stóp. Biodra znajdują się blisko powierzchni wody. Uniesiona głowa obniża biodra i nogi, co zwiększa przekrój ciała. Zbyt mocno opuszczona głowa przesuwa środek ciężaru do przodu i może zaburzać rytm oddychania.

Drugim elementem wyścigu na basenie krótkim lub 25-metrowym jest odbicie od ściany. W fazie ślizgu zawodnik porusza się bez napędu. Jego prędkość stopniowo maleje. Pływak powinien rozpocząć pracę nóg i rąk w momencie, gdy prędkość ślizgu spada do poziomu prędkości pływania. Zbyt wczesne rozpoczęcie pływania marnuje impet odbicia. Zbyt późne powoduje spowolnienie poniżej prędkości pływania.

Cel pracy: ustalić, która z trzech badanych pozycji głowy i która z czterech długości ślizgu daje najkrótszy czas przepłynięcia 25 m stylem dowolnym.

Hipotezy:
\begin{enumerate}
\item Pozycja neutralna jest szybsza od pozycji z głową uniesioną.
\item Dłuższy ślizg skraca czas przepłynięcia dystansu.
\end{enumerate}

\section*{2. Metoda}
\subsection*{2.1. Badany i warunki}
Badanie przeprowadzono na jednej osobie, autorze pracy. Pomiary wykonano na lokalnym basenie w Krasnym Stawie, na dystansie 25 m, stylem dowolnym. Czas mierzył stoperem pomocnik, który nie znał celu poszczególnych prób. Zawodnik wykonywał próby w stałej kolejności zmieniającej się losowo, z przerwą 3 minuty między próbami, aby ograniczyć wpływ zmęczenia.

\subsection*{2.2. Badane warianty}
\textbf{Część I: pozycja głowy.} Przy ślizgu o stałej długości 5 m zawodnik pływał w trzech wariantach.

\begin{tabularx}{\textwidth}{@{}lX@{}}
\toprule
\textbf{Wariant} & \textbf{Opis} \\
\midrule
A & pozycja neutralna: wzrok skierowany w dno, głowa w linii kręgosłupa \\
B & głowa uniesiona: wzrok skierowany do przodu \\
C & głowa opuszczona: wzrok skierowany do stóp, broda przy klatce piersiowej \\
\bottomrule
\end{tabularx}

\vspace{0.3cm}
Każdy wariant wykonano 8 razy.

\textbf{Część II: długość ślizgu.} W wariancie A zawodnik zmieniał długość ślizgu po odbiciu od ściany: 3 m, 4 m, 5 m i 6 m. Każdą długość wykonano 8 razy. Długość oznaczono pływakami ustawionymi na torze.

\subsection*{2.3. Opracowanie wyników}
Dla każdego wariantu obliczono średni czas oraz zakres wyników. Błąd pomiaru stoperem oceniono na około $\pm$0,2 s.

\section*{3. Wyniki}
\subsection*{3.1. Pozycja głowy}
\begin{tabularx}{\textwidth}{@{}lXX@{}}
\toprule
\textbf{Wariant} & \textbf{Średni czas [s]} & \textbf{Zakres [s]} \\
\midrule
A (neutralna) & 18,6 & 18,2--19,1 \\
B (głowa uniesiona) & 19,7 & 19,2--20,3 \\
C (głowa opuszczona) & 19,0 & 18,5--19,6 \\
\bottomrule
\end{tabularx}

\vspace{0.3cm}
Wariant A był szybszy od wariantu B o 1,1 s i od wariantu C o 0,4 s. Różnica między A i C jest dwukrotnie większa od oszacowanego błędu pomiaru. Różnica między A i B jest ponad pięciokrotnie większa.

\subsection*{3.2. Długość ślizgu}
\begin{tabularx}{\textwidth}{@{}lXX@{}}
\toprule
\textbf{Długość ślizgu} & \textbf{Średni czas [s]} & \textbf{Zakres [s]} \\
\midrule
3 m & 18,5 & 18,1--19,0 \\
4 m & 18,3 & 17,9--18,8 \\
5 m & 18,6 & 18,2--19,1 \\
6 m & 19,1 & 18,6--19,7 \\
\bottomrule
\end{tabularx}

\vspace{0.3cm}
Najkrótszy średni czas uzyskano przy ślizgu 4 m. Ślizg 6 m dał czas dłuższy o 0,8 s. Różnica między ślizgiem 3 m i 4 m mieści się w granicach błędu pomiaru.

\section*{4. Omówienie wyników}
Hipoteza pierwsza została potwierdzona. Pozycja neutralna była szybsza od pozycji z głową uniesioną. Wynik jest zgodny z zasadami hydrodynamiki. Uniesiona głowa obniża biodra i zwiększa opór kształtu. Zawodnik przed badaniem stosował właśnie wariant B. Pomiary pokazały więc, że jego nawykowa technika była najwolniejsza z badanych.

Hipoteza druga nie została potwierdzona. Dłuższy ślizg nie skracał czasu. Zależność miała kształt krzywej z minimum przy 4 m. Wyjaśnienie jest następujące. Prędkość w ślizgu maleje z czasem. Po przekroczeniu pewnej odległości zawodnik porusza się wolniej niż płynąc. Dalsze oczekiwanie na rozpoczęcie ruchów powoduje straty. Optymalna długość ślizgu zależy od prędkości odbicia i oporu ciała danego zawodnika. Zawodnik o innej budowie lub sile odbicia może uzyskać inne optimum.

Wynik części II był zaskoczeniem, ponieważ przed badaniem zawodnik zakładał, że dłuższy ślizg jest zawsze korzystniejszy. Zysk z korekty ślizgu z 5 m do 4 m wyniósł 0,3 s, a zysk z korekty pozycji głowy z B do A wyniósł 1,1 s. Łącznie korekty te dały poprawę o około 1,4 s na 25 m. Dla dystansu krótkiego jest to różnica duża w stosunku do typowych różnic między zawodnikami w zawodach lokalnych.

\section*{5. Ograniczenia}
\begin{itemize}
\item Badanie wykonano na jednej osobie, więc wyniki nie pozwalają na uogólnienia.
\item Czas mierzono stoperem, a błąd pomiaru wynosi około 0,2 s.
\item Liczba prób w wariancie wynosiła 8.
\item Pomiary wykonano na jednym basenie i na jednym dystansie.
\item Nie mierzono liczby cykli ramion ani prędkości odbicia.
\end{itemize}

\section*{6. Wnioski}
\begin{enumerate}
\item Najkrótszy średni czas na 25 m stylem dowolnym uzyskano przy neutralnej pozycji głowy.
\item Pozycja z głową uniesioną była wolniejsza od neutralnej o 1,1 s.
\item Optymalna długość ślizgu po odbiciu wyniosła 4 m. Ślizg dłuższy wydłużał czas przepłynięcia.
\item Zawodnik powinien dobierać długość ślizgu indywidualnie, na podstawie własnych pomiarów.
\item Dalsze badania wymagają większej liczby badanych, pomiaru elektronicznego i nagrań wideo.
\end{enumerate}

\vspace{0.4cm}
\textit{Pracę zatwierdził dr Robert Żukawski.}